% Imaging tests. Continues the chapter "The prototype so far" (2.0-prototype.tex).
% Draft by Claude, shortened on 2026-10-28 after feedback, edit freely.
% Red text in the PDF is a placeholder. Search for \fillin and \future.
% All test pictures with notes and settings: iterations/v0-concept/camera-settings-2026-10-26/
% Full-size originals: GitHub release data-2026-10-26.

\section{Imaging tests}\label{sec:imaging-tests}

The camera is the Raspberry Pi HQ camera with the lens removed. The tape lies above the
sensor and the light shines through it, so the picture is a shadow of what is on the tape.
One pixel is 1.55~$\mu$m, so a birch pollen grain of about 22~$\mu$m is about 14 pixels across.
As a test object we used a light brown hair on the tape (50 to 100~$\mu$m wide), until we can
get real pollen. \future{Repeat the tests with real pollen.}

The camera settings that worked best changed with every change to the setup. They will change
again when the camera sits in the box, so we do not report them in detail. The web interface
saves the settings with every picture, so we can go back to them
(Section~\ref{sec:control-software}).

\subsection{Test 1: tape 11.5~mm above the sensor}\label{sec:imaging-test-1}

In the first test the tape lay on top of the lens mount. Because of the IR filter in the mount
and the walls of the mount, the tape was about 11.5~mm above the sensor. We tried room light and
a blue LED about 30~cm above the tape (Figure~\ref{fig:camtest-best}).

The hair was visible, but blurred. Room light gave clearly better pictures than the blue LED.
With pure blue light only one in four pixels sees anything, because of the colour filter on the
sensor, and the pictures were dark and noisy. Because of this we removed the IR filter to make
the gap smaller, and replaced the blue LED with an RGB LED (Section~\ref{sec:electronics}).

\begin{figure}[htbp]
  \centering
  \begin{minipage}[t]{0.48\textwidth}
    \centering
    \includegraphics[width=\linewidth]{camtest-09}\\[2pt]
    \small (a) Room light
  \end{minipage}\hfill
  \begin{minipage}[t]{0.48\textwidth}
    \centering
    \includegraphics[width=\linewidth]{camtest-21}\\[2pt]
    \small (b) Blue LED
  \end{minipage}
  \caption{Test 1, tape 11.5~mm above the sensor. The dark line is the hair.}
  \label{fig:camtest-best}
\end{figure}

\subsection{Test 2: tape about 1~mm above the sensor}\label{sec:imaging-test-2}

For the second test the tape lies on the IR filter, which now rests about 1~mm above the sensor.
The filter is only there as a dust cover, so we can handle the setup without fear of damaging the
sensor. \fillin{Check the gap.}

With room light the hair is now much sharper than in Test~1 (Figure~\ref{fig:hairtest2}).
With the RGB LED, a single colour gave a sharper shadow than white. The three colours sit in
slightly different places inside the LED, so white light casts three shadows that are shifted a
little, which blurs the edges. Green and blue worked best.

The LED pictures are still more blurred than they should be. The LED is not a small point of light,
so each point on the tape casts a soft shadow. A pinhole in front of the LED, or a longer black tube
between LED and tape, should give a cleaner shadow. Putting the camera and LED inside the closed box
should also help, because no room light will reach the sensor.
\future{Add a pinhole or a black tube in front of the LED and repeat the test.}
\future{Repeat the test with the camera and LED inside the box, and without the IR filter.}

\begin{figure}[htbp]
  \centering
  \begin{minipage}[t]{0.48\textwidth}
    \centering
    \includegraphics[width=\linewidth]{hairtest2-roomlight}\\[2pt]
    \small (a) Room light
  \end{minipage}\hfill
  \begin{minipage}[t]{0.48\textwidth}
    \centering
    \fbox{\parbox[b][0.72\linewidth][c]{0.96\linewidth}{\centering\future{Picture with the best LED colour, after the pinhole or tube}}}\\[2pt]
    \small (b) Best LED colour
  \end{minipage}
  \caption{Test 2, tape about 1~mm above the sensor.}
  \label{fig:hairtest2}
\end{figure}
